\begin{table}[htbp]
\centering
\small
\caption{Bloco 2, Problema E2 (openai/gpt-4o-mini): aprendizado em contexto por estratégia de seleção e n.}
\label{tab:bloco2_E2_faseE_estrategia_x_nshot_gpt4mini}
\begin{tabular}{lrrrrrrrr}
\toprule
 & \multicolumn{2}{c}{Fácil} & \multicolumn{2}{c}{Difícil} & \multicolumn{2}{c}{Mista} & \multicolumn{2}{c}{Aleatória} \\ \cmidrule(lr){2-3} \cmidrule(lr){4-5} \cmidrule(lr){6-7} \cmidrule(lr){8-9}
n & Acur.\ vs espec. & $\kappa$ & Acur.\ vs espec. & $\kappa$ & Acur.\ vs espec. & $\kappa$ & Acur.\ vs espec. & $\kappa$ \\
\midrule
zero-shot (único valor) & 48,9 (46,7–51,3) & 0,002 &  &  &  &  &  &  \\
4-shot & 73,3 (69,2–76,7) & 0,476 & 41,1 (31,5–48,6) & -0,177 & 79,9 (79,5–80,8) & 0,599 & 81,8 (54,1–94,5) & 0,639 \\
10-shot & 79,0 (75,0–82,1) & 0,588 & 79,3 (71,4–84,3) & 0,579 & 87,1 (77,9–97,1) & 0,742 & 87,7 (77,9–94,3) & 0,755 \\
20-shot & 80,0 (75,4–83,1) & 0,607 & 95,9 (93,1–100,0) & 0,918 & 90,8 (85,4–98,5) & 0,816 & 94,1 (86,9–96,9) & 0,882 \\
40-shot & 77,6 (75,5–79,1) & 0,562 & 92,7 (87,3–100,0) & 0,856 & 94,5 (90,9–99,1) & 0,891 & 96,0 (91,8–98,2) & 0,919 \\
\bottomrule
\multicolumn{9}{p{0.95\linewidth}}{\footnotesize Especialista E2, W = [1,5; 0,3]: rotula o Problema E pelo centróide mais próximo; o LLM recebe n exemplos escolhidos pela estratégia e classifica os demais. Acur.\ vs espec.\ = acurácia (LLM vs especialista): fração dos pontos de teste em que o LLM reproduz o rótulo do especialista. Valores em \% (mín–máx); $\kappa$ = Cohen's $\kappa$ (faixas de Landis e Koch, 1977). Médias das sementes e sorteios. Em zero-shot não há exemplo a selecionar: valor único.} \\
\end{tabular}
\end{table}
